\begin{frame}[t]{Perspective : adapter le poids de la hiérarchie à l'image}
  \vspace{0.06cm}
  \begin{columns}[T,onlytextwidth]
    \begin{column}{0.35\textwidth}
      \small\textbf{Une motivation : les ombres}\par\medskip
      \begin{minipage}[t]{0.48\linewidth}\centering
        \includegraphics[width=\linewidth]{ombre_image.jpg}\par
        {\scriptsize Image}
      \end{minipage}\hfill
      \begin{minipage}[t]{0.48\linewidth}\centering
        \includegraphics[width=\linewidth]{ombre_frangi.jpg}\par
        {\scriptsize Réponse Frangi}
      \end{minipage}\par\smallskip
      {\scriptsize Exemple de nos essais : l'ombre attire Frangi.}
    \end{column}
    \begin{column}{0.60\textwidth}
      \small\textbf{Faire prédire un poids par les caractéristiques SAM}\par\medskip
      \centering
      \resizebox{0.96\linewidth}{!}{% Proposition nouvelle : coefficient global appris, aucune mesure de performance.
\begin{tikzpicture}[every node/.style={font=\small,align=center},
  confidencebox/.style={draw=gris_fonce_inria,rounded corners=3pt,fill=gray!4,
    minimum height=.85cm,inner sep=5pt},
  flow/.style={-{Latex[length=2.4mm]},draw=gris_fonce_inria,line width=.9pt}]
  \node[confidencebox] (features) at (0,0) {Caractéristiques\\SAM : $\bar h(I)$};
  \node[confidencebox,draw=inria-rouge] (head) at (3.5,0) {Petite couche\\apprise};
  \node[confidencebox,draw=inria-rouge,fill=inria-rouge!5] (beta) at (6.65,0) {$\beta(I)$\\entre 0 et 1};
  \draw[flow] (features)--(head); \draw[flow] (head)--(beta);
\end{tikzpicture}
}
      \[\beta(I)=\mathrm{sigmoid}\!\left(w^\top\bar h(I)+b\right).\]
      {\scriptsize $\bar h(I)$ : moyenne des caractéristiques \textbf{avant} l'attention guidée.}
    \end{column}
  \end{columns}
  \vspace{0.16cm}
  \begin{center}
    \formulebox{\parbox{0.81\textwidth}{\centering\small
      $A_H(I)=\mathrm{softmax}\!\left(\text{scores visuels}+\beta(I)B_H(I)\right)$}}
  \end{center}
  \vspace{0.08cm}
  \small
  \begin{itemize}
    \setlength{\itemsep}{3pt}
    \item \textbf{Hypothèse :} moins de guidage quand Frangi gêne, notamment dans les ombres
      ou la granularité. Poids appris avec LoRA par la perte de segmentation.
    \item \textbf{Test :} sans biais / $\beta$ constant / $\beta(I)$.
      Comparer IoU, ruptures et faux raccords sur les mêmes images.
  \end{itemize}
  \vspace{0.06cm}
  {\footnotesize $B_H(I)$ dépend déjà de l'image ; on adapte ici \textbf{son poids}.
  Un poids global peut aussi réduire le guidage dans des zones utiles.}
\end{frame}
